% Standalone, self-contained Overleaf document. Compile with pdfLaTeX.
% To reuse in a manuscript, copy Section 1 and the required macros/packages.
% This describes the agreed conceptual framework, not every classifier exception.
\documentclass[11pt]{article}
\usepackage[T1]{fontenc}
\usepackage{lmodern}
\usepackage[margin=1in]{geometry}
\usepackage{amsmath,amssymb}
\usepackage{booktabs,tabularx,array}
\usepackage{microtype}
\usepackage[hidelinks]{hyperref}

\newcommand{\AIH}{\mathrm{AI}\mid\mathrm{H}}
\newcommand{\RPE}{\mathrm{RPE}}
\newcommand{\SPE}{\mathrm{SPE}}
\newcommand{\At}{\mathcal{A}_t^{\AIH}}
\newcommand{\Et}{\mathcal{E}_t^{\AIH}}
\newcommand{\gt}{g_t^{\AIH}}
\newcommand{\ind}{\mathbf{1}}
\newcolumntype{Y}{>{\raggedright\arraybackslash}X}
\setlength{\emergencystretch}{2em}

\title{Stepwise Reward and State Prediction Errors\\in Human--AI Collaboration}
\author{}
\date{Working Methods formulation --- 16 October 2026}

\begin{document}
\maketitle

\section{Methods}
\label{sec:methods}

\subsection{Collaborative state and human expectations}
We describe collaboration using a goal-conditioned reinforcement-learning
notation. At step $t$, $s_t$ is the shared state, including the task context and
relevant interaction history; $a_t^{\mathrm H}$ and $a_t^{\mathrm{AI}}$ are the
human and AI actions. The environment has a transition model $T$, and
$s_{t+1}$ denotes the observed state after the joint action. The task the human
expects the AI to pursue is $\gt$, where the superscript $\AIH$ means
``expected of the AI by the human.'' This expectation is established from
instructions, reviewed demonstrations, and the current task phase; an
erroneous AI action does not itself redefine it.

Following the distinction between goal-related and state-transition-related
discrepancies in neural value alignment~\cite{xu2026nva}, we assign two binary
labels to each step. Reward prediction error (RPE) indicates incompatibility
with the expected task. State prediction error (SPE) indicates that the
observed next state falls outside acceptable execution of that task. These
are operational behavioral labels, rather than estimates of a participant's
neural response or numerical temporal-difference errors.

\subsection{RPE: compatibility with the expected task}
We define a binary task-alignment reward,
\begin{equation}
 r_{t+1}^{G}
 =R_G^{\AIH}\!\left(s_t,a_t^{\mathrm H},a_t^{\mathrm{AI}},s_{t+1};\gt\right)
 \in\{0,1\}.
 \label{eq:reward}
\end{equation}
Here, $G$ identifies the goal-alignment component, and $R_G^{\AIH}$ evaluates
whether the AI contribution supports the human's expected task.
A value of $1$ denotes compatible progress or an established justification
for preparation, waiting, coordination, or recovery. A value of $0$ denotes
a contradiction of the required task, content, recipient, or destination,
or unjustified regression or nonprogress when a feasible advancing action
exists. These criteria are operationalized for each domain using the
expectation evidence and current constraints. Immediate task completion is
not required at every step.

Taking compatible behavior, $r_{t+1}^{G}=1$, as the normative reference gives
\begin{equation}
 \boxed{\RPE_t=1-r_{t+1}^{G}.}
 \label{eq:rpe}
\end{equation}
Thus, $\RPE_t=1$ means that the observed contribution is incompatible with
the expected task; $\RPE_t=0$ means that it is compatible. This reward is
an annotation of task alignment, not the environment's game score, a
learned reward model, or an empirically measured human reward expectation.

\subsection{SPE: compatibility with acceptable execution}
Let $\At$ be the set of AI actions accepted under the established human
preferences at the current state and task phase, conditional on
$a_t^{\mathrm H}$. This set is specified directly from task instructions,
reviewed demonstrations, and coordination constraints; it need not be
obtained from a probabilistic policy. Multiple actions can be acceptable.

Let $\phi$ map a state to the next-state features used to evaluate execution.
Unlike the scalar reward in Equation~\eqref{eq:reward}, $\phi(s_{t+1})$
describes what happened, such as changes in agent positions, shared objects,
or a jointly edited document. It preserves distinctions that a single
task-alignment score would discard. For a deterministic transition model
conditional on the specified actions, the accepted next-state
representations are
\begin{equation}
 \Et
 =\left\{\phi\!\left(T(s_t,a_t^{\mathrm H},a)\right)
          :a\in\At\right\}.
 \label{eq:expected}
\end{equation}
The symbol $a$ is a candidate AI action. The braces collect the next-state
representation produced by each accepted candidate while keeping the same
starting state and human action. SPE is then
\begin{equation}
 \boxed{\SPE_t=\ind\!\left[\phi(s_{t+1})\notin\Et\right].}
 \label{eq:spe}
\end{equation}
The indicator $\ind[\cdot]$ equals $1$ when its condition is true and $0$
otherwise. Therefore, execution is accepted if its observed representation
matches at least one accepted alternative. Different commands can receive
the same label when they produce the same relevant state transition.

The two labels answer different questions: whether the contribution
supports the expected task, and whether its execution is acceptable.
They may rely on overlapping evidence and may both equal $1$; they are not
assumed to be statistically independent. When the expected task or accepted
execution cannot be established, the corresponding label remains
unresolved. The fixed-human-action comparison is a retrospective,
conditional assessment, not a requirement that the AI anticipate an
unobserved simultaneous human action.

\subsection{Operationalization in Overcooked}
In Overcooked, the shared state contains both chefs' positions,
orientations and held items, counter objects, pot contents and cooking
states, remaining orders, and the established task context. Each chef can
move north, south, east or west, stay, or interact. Expected AI subtasks
include acquiring an onion, tomato or plate; receiving an item supplied by
the human; supplying an item to the human through a handoff counter;
adding an ingredient; plating soup; and serving. Receiving and supplying
are distinct directional roles.

The current subtask follows reviewed action blocks and prerequisite
completion. For example, acquiring a plate remains the immediate task
until pickup; supplying it begins on the following action. An interaction
that successfully completes a subtask belongs to that subtask's block.
If the AI deviates, the human's preferred task remains the reference.
A particular counter, source or pot is a goal constraint when the Hint or
reviewed context requires it. Otherwise, targets that satisfy the subtask
can remain goal-compatible, although execution preferences can still
favor a particular target or route.

For RPE, appropriate approach, necessary facing, human-blocked intended
progress, and justified detours or yielding can receive $r_{t+1}^{G}=1$.
Consequently, unchanged distance alone does not establish a goal error.
For SPE, accepted actions respect the preferred target and established
human passage, including justified waiting and clearance. The production
transition model evaluates candidate AI commands from the same Before
state with the recorded human command held fixed. The representation
$\phi$ includes both chefs' poses and inventories, world objects and pot
states, and remaining orders; it excludes animation and game score.

For example, suppose the human expects an onion handoff through an accepted
counter while retaining access to an established passage.
Table~\ref{tab:handoff} illustrates the resulting distinctions. A
goal-consistent approach can still produce SPE when it unnecessarily
obstructs that passage. Conversely, moving temporarily away to yield can
be compatible with both expectations.

\begin{table}[htbp]
 \centering\small
 \caption{Illustrative labels under an established onion-handoff expectation.}
 \label{tab:handoff}
 \begin{tabularx}{\linewidth}{@{}Ycc@{}}
 \toprule
 AI contribution & $\RPE_t$ & $\SPE_t$ \\
 \midrule
 Accepted approach or handoff & 0 & 0 \\
 Pursues the handoff but unnecessarily blocks the human's passage & 0 & 1 \\
 Adds the onion to a pot instead of supplying the human & 1 & 1 \\
 Makes an accepted yield or justified wait to enable passage & 0 & 0 \\
 \bottomrule
 \end{tabularx}
\end{table}

The equations summarize the conceptual method; the current audit further
operationalizes readiness, recovery, reference matching, and coordination
through ordered rules. In particular, a reviewed incompatible-input state
can establish SPE even when an executable recovery alternative is not
available. Computation is restricted to saved versions with current
Playback and Subtasks checks and a correspondingly checked nonPE
reference. Authored condition names and critical/correction markers do not
determine individual step labels.

\subsection{Application beyond physical collaboration}
For LLM-assisted writing, $s_t$ can contain the shared document,
instructions and relevant conversation; an AI action can be an entire
response, edit or tool call, and $T$ updates that shared state. If the
expected task is to summarize Article A in three cited bullets, a faithful
summary in the accepted format can receive $(\RPE_t,\SPE_t)=(0,0)$.
A faithful summary in an unrequested prose format can receive $(0,1)$,
whereas translating instead of summarizing can receive $(1,1)$.
Requesting genuinely missing source text can be accepted preparation.
Here, $\phi$ describes execution properties such as structure, citation
preservation and document modifications; explicit annotation or a
validated semantic evaluator replaces exact kitchen-state comparison.
These examples illustrate a proposed transfer, not an implemented LLM
classifier. When next-state features remain stochastic conditional on
specified actions, Equation~\eqref{eq:expected} requires a calibrated
predictive acceptance region rather than exact point matching.

\appendix
\section{Notation and reading the equations}
\label{sec:notation}
\begin{table}[htbp]
 \centering\small
 \begin{tabularx}{\linewidth}{@{}lY@{}}
 \toprule
 Symbol & Meaning \\
 \midrule
 $t$ & Index of the joint action being evaluated. \\
 $s_t,\ s_{t+1}$ & Shared state before and after that joint action. \\
 $a_t^{\mathrm H},\ a_t^{\mathrm{AI}}$ & Observed human and AI actions. \\
 $T$ & Environment transition model, deterministic conditional on actions
 in Equation~\eqref{eq:expected}. \\
 $\AIH$ & Expected of the AI by the human. \\
 $\gt$ & Human's currently established preferred AI task. \\
 $R_G^{\AIH},\ r_{t+1}^{G}$ & Task-alignment evaluator and its binary value;
 $G$ denotes goal alignment. \\
 $\At$ & Set of acceptable candidate AI actions in the current context. \\
 $a$ & One candidate AI action used in the comparison. \\
 $\phi$ & Mapping from a state to execution-relevant state features. \\
 $\Et$ & Set of next-state representations generated by accepted actions. \\
 $\ind[\cdot]$ & Indicator: $1$ for a true condition, otherwise $0$. \\
 $\notin$ & Is not an element of the specified set. \\
 $\RPE_t,\ \SPE_t$ & Binary task-mismatch and execution-mismatch labels. \\
 \bottomrule
 \end{tabularx}
\end{table}

\paragraph{Reading the set notation.}
Suppose that, at one particular state, both moving east and staying are
accepted. Then $\At=\{\mathrm{EAST},\mathrm{STAY}\}$ and
Equation~\eqref{eq:expected} expands to
\begin{equation*}
 \begin{split}
 \Et=\bigl\{&\phi(T(s_t,a_t^{\mathrm H},\mathrm{EAST})),\\
             &\phi(T(s_t,a_t^{\mathrm H},\mathrm{STAY}))\bigr\}.
 \end{split}
\end{equation*}
The colon in Equation~\eqref{eq:expected} means ``for each accepted action.''
There is no policy probability or probability threshold in this
definition. If the observed next-state representation equals either
listed representation, SPE is $0$; otherwise it is $1$. This is an
illustration of the notation, not a claim that east and stay are always
acceptable together.

\paragraph{Why state features are different from reward.}
The binary reward compresses the contribution into a task-compatibility
judgment. The state representation retains execution details. Two
goal-compatible contributions can therefore both earn reward $1$ while
differing in whether they obstruct the human or preserve a requested
document structure. The feature mapping should be specified before
classification: retaining only reward would discard precisely the
execution differences that SPE is intended to detect.

\begin{thebibliography}{1}
\bibitem{xu2026nva}
X.~Xu, Y.~Wang, D.~Han, D.~Li, and S.~W.~Lee,
``Neural Value Alignment: Human--AI Collaboration Under Goal-Action Ambiguity,''
\emph{IEEE Transactions on Cybernetics}, 2026.
\href{https://doi.org/10.1109/TCYB.2026.3722605}{doi:10.1109/TCYB.2026.3722605}.
\end{thebibliography}

\end{document}
